\documentclass[10pt,a4paper]{amsart}
\usepackage{amsmath,amssymb,amsthm,graphicx}
\usepackage{verbatim}
\newenvironment{reference}{\comment}{\endcomment}
%\newenvironment{reference}{}{}
\newenvironment{slogan}{\comment}{\endcomment}
\newenvironment{history}{\comment}{\endcomment}
% For commutative diagrams we use Xy-pic
\usepackage[all]{xy}

% We use 2cell for 2-commutative diagrams.
\xyoption{2cell}
% Theorem environments.
%
\theoremstyle{plain}
\newtheorem{theorem}[subsection]{Theorem}
\newtheorem{proposition}[subsection]{Proposition}
\newtheorem{lemma}[subsection]{Lemma}

\theoremstyle{definition}
\newtheorem{definition}[subsection]{Definition}
\newtheorem{example}[subsection]{Example}
\newtheorem{exercise}[subsection]{Exercise}
\newtheorem{situation}[subsection]{Situation}

\theoremstyle{remark}
\newtheorem{remark}[subsection]{Remark}
\newtheorem{remarks}[subsection]{Remarks}

\numberwithin{equation}{subsection}

% Macros
%
\def\lim{\mathop{\mathrm{lim}}\nolimits}
\def\colim{\mathop{\mathrm{colim}}\nolimits}
\def\Spec{\mathop{\mathrm{Spec}}}
\def\Hom{\mathop{\mathrm{Hom}}\nolimits}
\def\Ext{\mathop{\mathrm{Ext}}\nolimits}
\def\SheafHom{\mathop{\mathcal{H}\!\mathit{om}}\nolimits}
\def\SheafExt{\mathop{\mathcal{E}\!\mathit{xt}}\nolimits}
\def\Sch{\mathit{Sch}}
\def\Mor{\mathop{\mathrm{Mor}}\nolimits}
\def\Ob{\mathop{\mathrm{Ob}}\nolimits}
\def\Sh{\mathop{\mathit{Sh}}\nolimits}
\def\NL{\mathop{N\!L}\nolimits}
\def\CH{\mathop{\mathrm{CH}}\nolimits}
\def\proetale{{pro\text{-}\acute{e}tale}}
\def\etale{{\acute{e}tale}}
\def\QCoh{\mathit{QCoh}}
\def\Ker{\mathop{\mathrm{Ker}}}
\def\Im{\mathop{\mathrm{Im}}}
\def\Coker{\mathop{\mathrm{Coker}}}
\def\Coim{\mathop{\mathrm{Coim}}}

% Boxtimes
%
\DeclareMathSymbol{\boxtimes}{\mathbin}{AMSa}{"02}

%
% Macros for moduli stacks/spaces
%
\def\QCohstack{\mathcal{QC}\!\mathit{oh}}
\def\Cohstack{\mathcal{C}\!\mathit{oh}}
\def\Spacesstack{\mathcal{S}\!\mathit{paces}}
\def\Quotfunctor{\mathrm{Quot}}
\def\Hilbfunctor{\mathrm{Hilb}}
\def\Curvesstack{\mathcal{C}\!\mathit{urves}}
\def\Polarizedstack{\mathcal{P}\!\mathit{olarized}}
\def\Complexesstack{\mathcal{C}\!\mathit{omplexes}}
% \Pic is the operator that assigns to X its picard group, usage \Pic(X)
% \Picardstack_{X/B} denotes the Picard stack of X over B
% \Picardfunctor_{X/B} denotes the Picard functor of X over B
\def\Pic{\mathop{\mathrm{Pic}}\nolimits}
\def\Picardstack{\mathcal{P}\!\mathit{ic}}
\def\Picardfunctor{\mathrm{Pic}}
\def\Deformationcategory{\mathcal{D}\!\mathit{ef}}
\usepackage{cleveref}
\crefname{originalreference}{original reference}{original references}
\Crefname{originalreference}{Original reference}{Original references}
\makeatletter
\newlabel{dualizing-lemma-torsion-change-rings}{{dualizing-lemma-torsion-change-rings}{source}{}{sourcepointer-1}{}}
\newlabel{dualizing-lemma-torsion-change-rings@cref}{{[originalreference][0][]dualizing-lemma-torsion-change-rings}{[1][0][]source}{}{}{}}
\makeatother
\title{Base change of the sections with support along a flat ring map, and the bounded case when the rings have finite dimension}
\author{The Stacks Project Authors}
\date{stacks-project:local-cohomology.tex:local-snapshot-20261012; GNU FDL 1.2 or later}
\begin{document}
\maketitle
\noindent\textit{Editorial note.} All delivered mathematical text is the author's own, copied from the pinned Stacks Project chapter source sources/stacks-local-cohomology/pinned-original-source.tex; no statement or proof step was rewritten, reordered or removed.\par\medskip
\noindent\textit{Editorial note.} Preamble provenance: the Stacks Project keeps one single shared preamble for every chapter; this card takes it directly from the staged sources/stacks-local-cohomology/preamble.tex as three verbatim spans (lines 8-12, 14-18 and 154-217), so no derived file was written and nothing was added or redefined. The wrapper class and the base packages are supplied by the assembler.\par\medskip
\noindent\textit{Editorial note.} Bibliography provenance: neither the counted statement nor its proof contains a citation (this lane rescans the whole span before writing the card), so the card carries no bibliography and no entry was reconstructed.\par\medskip
\noindent\textit{Editorial note.} Selection provenance: this card was opened under the file-based control contract (CONTROL.md), with the mechanical pre-card gate run on its statement and proof interval before the card was written.\par\medskip
\noindent\textit{Allowed standard prerequisites (editorial).}\begin{itemize}
\item The description of the local cohomology with support in a specialization-stable subset as the filtered colimit of the local cohomologies with support in the closed subsets contained in it. Author lemma lemma-adjoint-ext, whose statement is reproduced verbatim as a context span of this card (source lines 996-1004). Reduces the general support to the closed case by passing to the colimit over the closed subsets, on both sides of the base change map.
\item The base change statement for the sections with support in a closed subset along a flat map, in the form recorded in the chapter on dualizing complexes. Invoked inside the counted proof at source lines 1118-1119 as a cross-chapter pointer. Provides the closed case from which the general case is assembled.
\end{itemize}\medskip
\ifdefined\hypertarget\hypertarget{sourcepointer-1}{}\fi
\noindent\textit{Source reference.} Original reference dualizing-lemma-torsion-change-rings: sources/stacks-local-cohomology/pinned-original-source.tex, source lines 741--741. Navigation pointer only: the counted text refers to the Stacks item dualizing-lemma-torsion-change-rings, which belongs to another chapter of the project and is not part of this card. The pointer carries the target's own label as its number because this chapter's staged snapshot does not contain the other chapter's tag table; its content is not reproduced here.\par\medskip
\begin{lemma}
\label{lemma-torsion-change-rings}
Let $A \to B$ be a flat homomorphism of Noetherian rings.
Let $T \subset \Spec(A)$ be a subset stable under specialization.
Let $T' \subset \Spec(B)$ be the inverse image of $T$.
Then the canonical map
$$
R\Gamma_T(K) \otimes_A^\mathbf{L} B
\longrightarrow
R\Gamma_{T'}(K \otimes_A^\mathbf{L} B)
$$
is an isomorphism for $K \in D^+(A)$. If $A$ and $B$ have finite
dimension, then this is true for $K \in D(A)$.
\end{lemma}

\begin{lemma}
\label{lemma-adjoint-ext}
Let $A$ be a Noetherian ring. Let $T \subset \Spec(A)$
be a subset stable under specialization.
For any object $K$ of $D(A)$ we have
$$
H^i(RH^0_T(K)) = \colim_{Z \subset T\text{ closed}} H^i_Z(K)
$$
\end{lemma}

\begin{proof}
From the map $R\Gamma_T(K) \to K$ we get a map
$R\Gamma_T(K) \otimes_A^\mathbf{L} B \to K \otimes_A^\mathbf{L} B$.
The cohomology modules of $R\Gamma_T(K) \otimes_A^\mathbf{L} B$
are supported on $T'$ and hence we get the arrow of the lemma.
This arrow is an isomorphism if $T$ is a closed subset of $\Spec(A)$ by
Dualizing Complexes, Lemma \ref{dualizing-lemma-torsion-change-rings}.
Recall that $H^i_T(K)$ is the colimit of $H^i_Z(K)$ where $Z$ runs over
the (directed set of) closed subsets of $T$, see
Lemma \ref{lemma-adjoint-ext}.
Correspondingly
$H^i_{T'}(K \otimes_A^\mathbf{L} B) =
\colim H^i_{Z'}(K \otimes_A^\mathbf{L} B)$ where $Z'$ is the inverse
image of $Z$. Thus the result because $\otimes_A B$ commutes
with filtered colimits and there are no higher Tors.
\end{proof}

\end{document}
